% SCAFFOLD B: "Subspaces encode entity bindings, not beliefs"
% Use if: E14 false-belief IIA < 0.30, E12 shows framing-invariance,
%          E15 observability IIA < 0.30
% Current data: UNLIKELY given E14 L38=0.975, but scaffold for completeness.
%
% This file contains ONLY the new/changed sections.
% The existing paper (lookback_audit.tex) stays intact through Section 5.
% Changes: new abstract, new Section 6 (Results), rewritten Section 7 (Discussion).

%%% ---- UPDATED ABSTRACT ---- %%%

% Mechanistic interpretability routinely labels causal findings with
% vocabulary borrowed from cognitive science without testing whether the
% evidence supports the specificity these labels imply. We audit a
% three-part ``lookback mechanism'' reported to track characters' beliefs
% in Llama-3.1-70B-Instruct through pointer dereference implemented via
% attention. A pre-registered validity assessment identifies six untested
% assumptions, most critically that the paper's evaluation never tests
% false beliefs: the story-generation code includes belief = reality for
% every evaluated pair. We conduct 15 pre-registered experiments. The
% subspaces survive shuffled-label and random-baseline controls,
% confirming they capture genuine task-relevant structure. However, they
% are framing-invariant (belief, action-recall, and reality questions
% produce matched IIA), fail to generalize to genuine false-belief
% stories (IIA < 0.30 when belief $\neq$ reality), and show no transfer
% to observability-mediated knowledge. We recharacterize the mechanism
% as attention-mediated entity-state retrieval: the subspaces encode
% ``what did this character put where?'' rather than ``what does this
% character believe?'' The mechanism is real; the label overshoots.

%%% ---- SECTION 6: RESULTS ---- %%%

\section{Results}
\label{sec:results}

\subsection{Positive controls and measurement validity (E1--E4)}

% Same as Scaffold A — the positive controls and measurement validity
% experiments don't depend on the belief-vs-binding question.
% E1: positive control replicates
% E2: shuffled labels beaten
% E3: task specificity passes (echo=0, entity_tracking=0)
% E4: surface heuristic results

\subsection{Internal structure (E5--E9)}

% Same structure as Scaffold A — these are architectural tests.

\subsection{Specificity of the belief interpretation (E10--E11)}

% E11 cross-task: L38 = 0.529 on factual recall becomes MORE important
% in this version — it's evidence that L38 is a generic retrieval
% mechanism, not belief-specific.

\subsection{Belief-vs-binding dissociation (E12--E15)}
\label{sec:dissociation}

% E12: Question-framing control
% - All framings produce matched IIA
% - INTERPRETATION: The subspace encodes entity-state bindings.
%   "What does X believe Y contains?" and "What did X put in Y?"
%   and "What is inside Y?" all activate the same subspace equally.
%   The belief framing adds no explanatory power.

% E13: Distractor insertion
% - If IIA drops with distractors: positional features
% - If IIA preserved: semantic binding (supports entity binding)

% E14: True false-belief test  *** KEY EXPERIMENT ***
% - Behavioral accuracy: low (< 0.30) OR
%   Behavioral accuracy > 0.50 but IIA < 0.30
% - INTERPRETATION (if behavioral accuracy > 0.50 AND IIA < 0.30):
%   "The subspace trained on entity-binding does not generalize to
%   false beliefs. Combined with E12 showing framing-invariance,
%   this supports the recharacterization: the subspace encodes
%   character-state bindings that happen to align with beliefs in
%   the paper's task, but does not encode epistemic state per se."

% E15: Observability dissociation  *** KEY EXPERIMENT ***
% - Observed-other-container IIA < 0.30
% - Self-container control matches baseline
% - INTERPRETATION: "The subspace only encodes self-action bindings.
%   It tracks 'what did I put where?' but not 'what did I see
%   someone else put where?' — the latter requires modeling
%   observability, which the subspace does not capture."

% ---- KEY SYNTHESIS ----
% 1. Framing-invariant (E12): not belief-specific
% 2. Fails on false beliefs (E14): not epistemic
% 3. Fails on observability (E15): not tracking who-sees-what
% 4. Partially transfers to factual recall (E11): generic retrieval
%
% The subspace encodes "what entity is associated with what state
% at what location" — entity-state binding. In the paper's task,
% this happens to align with belief because belief = reality.
% The "belief tracking" label overshoots the mechanism's scope.

\subsection{Summary table}

% TABLE: experiment | subspace | prediction | result | verdict
% E1-E4: green (measurement validity confirmed)
% E5-E9: mixed
% E10-E11: yellow (L38 partially generic)
% E12: red for belief interpretation (framing-invariant)
% E13: green or yellow
% E14: red (fails on false beliefs)
% E15: red (fails on observability)


%%% ---- SECTION 7: DISCUSSION (REWRITTEN) ---- %%%

\section{Discussion}
\label{sec:discussion_v2}

% PARAGRAPH 1: State the main finding
% The lookback subspaces are real — they survive every measurement
% control. The interpretive label is wrong — they do not track beliefs.

% PARAGRAPH 2: What the subspaces actually do
% The subspaces encode character-state associations: which character
% put which drink in which container. This is entity binding, a
% well-studied phenomenon in cognitive science and transformer
% interpretability. The paper's CausalToM task always queries
% characters about their own containers, so entity binding and belief
% produce identical answers. The false-belief test (E14) breaks this
% alignment and reveals the subspace tracks the binding, not the belief.

% PARAGRAPH 3: Why this matters
% The mechanism is genuinely interesting — it shows the model has
% structured representations for entity-state associations that can
% be isolated to specific layer ranges and low-rank subspaces. The
% contribution is real. The issue is referential inflation: calling
% this "belief tracking" imports cognitive specificity the evidence
% does not support. A more accurate characterization — "entity-state
% binding via attention-mediated retrieval" — preserves the technical
% contribution while avoiding overclaimed cognitive parallels.

% PARAGRAPH 4: Implications for mechanistic interpretability
% This case study illustrates a general pattern: interchange
% interventions identify WHAT information a subspace carries but not
% HOW the model uses it. A subspace achieving high IIA on a task
% demonstrates causal relevance to the task. But the task determines
% what the subspace is tested for — and a task where belief = reality
% cannot distinguish belief from entity binding. The interpretive
% label must be calibrated to the most parsimonious explanation
% consistent with the full experimental battery, not the most
% cognitively impressive one consistent with the original task.

% PARAGRAPH 5: Verdict
% We maintain the "Causally Suggestive" rating. The subspaces are
% causally relevant to the task (confirmed by E1, E2, E7). The
% mechanism label requires correction from "belief tracking" to
% "entity-state retrieval." The three-stage architecture (binding →
% answer → visibility) is better characterized as entity encoding →
% entity retrieval → input gating.

% PARAGRAPH 6: Limitations
% Same as Scaffold A

\paragraph{Contracted claim (revised).}
% "Belief tracking via pointer dereference" should be contracted to
% "entity-state retrieval via attention-mediated subspace routing."
% The pointer metaphor is apt — the subspace does carry a reference
% from the question to the relevant entity-state pair. But the
% referent is an entity-state binding, not an epistemic state.
